\chapter{COMSOL Parameter and Scenario Organization}

\section{Parameter Files}

The COMSOL-facing parameter files are stored under
\texttt{data/model\_parameters/}. The accelerated degradation study uses the
CSV parameter set for:
\begin{itemize}
\item master degradation state, \texttt{General Degradation D\_Master.csv};
\item scenario controls, \texttt{Case Study.csv};
\item series resistance, \texttt{Series Resistance Rs.csv};
\item shunt resistance, \texttt{Shunt Resistance Rsh.csv};
\item trap density, \texttt{Trap Density Nt (Bulk Recomb).csv};
\item mobile/interfacial charge, \texttt{Mobile Ions.csv}; and
\item surface recombination, \texttt{Surface Recomb S.csv}.
\end{itemize}

The accelerated amplitudes are intentionally strong enough to make the full
coupled case cross $T_{80}$ before 3000 h. This makes the library suitable for
mechanism sensitivity and calibration planning.

\section{Architecture Parameters}

The laboratory stack is
\begin{equation}
\mathrm{glass}/\mathrm{ITO}/\mathrm{NiO}_x/
\mathrm{MeO\mbox{-}2PACz}/\mathrm{MHP}/
\mathrm{C}_{60}/\mathrm{BCP}/\mathrm{Ag}/
\mathrm{UV\ resin}/\mathrm{glass}.
\end{equation}
The reduced electronic geometry uses the following model parameters:
\begin{itemize}
\item \texttt{thickness\_HTL = 35[nm]};
\item \texttt{thickness\_PVK = 450[nm]}; and
\item \texttt{thickness\_ETL = 29[nm]}.
\end{itemize}
The HTL combines approximately 30 nm NiO$_x$
and 5 nm MeO-2PACz; the ETL combines approximately 25 nm C$_{60}$ and 4 nm
BCP. The 100 nm Ag layer, ITO, resin, and glass are represented as contacts,
illumination context, or omitted encapsulation rather than semiconductor
transport domains.

\section{Scenario Semantics}

\begin{table}[H]
\centering
\begin{tabular}{lll}
\toprule
Mode & Case ID & Meaning \\
\midrule
0 & 0 & full coupled degradation \\
1 & 0 & all-off control \\
1 & 1--5 & one-only mechanism case \\
0 & 1--5 & leave-one-out mechanism case \\
\bottomrule
\end{tabular}
\caption{Scenario semantics used by the COMSOL sweep and processing pipeline.}
\label{tab:scenario_semantics}
\end{table}

For \texttt{case\_id=1..5}, the selected mechanisms are $R_s$, $R_{sh}$,
$N_t$, $Q_f$, and $S$.

\section{Generated Data Products}

The primary processed COMSOL products are:
\begin{itemize}
\item \texttt{data/processed/comsol/comsol\_scenario\_lookup.csv};
\item \texttt{data/processed/comsol/comsol\_scenario\_timeseries.csv};
\item \texttt{data/processed/comsol/comsol\_lifetime\_summary.csv};
\item \texttt{outputs/figures/comsol\_baseline\_pce.png};
\item \texttt{outputs/figures/comsol\_full\_jv\_evolution.png};
\item \texttt{outputs/figures/comsol\_mechanism\_sensitivity.png};
\item \texttt{outputs/figures/comsol\_metric\_signatures.png}; and
\item \texttt{outputs/figures/comsol\_state\_trajectories.png}.
\end{itemize}

These files are regenerated by
\texttt{scripts/run\_comsol\_paper\_pipeline.py}.

\section{Electronic Parameter-Sweep Products}

The separate electronic parameter sweep uses the COMSOL Global Evaluation data
under \path{data/raw/comsol/}:
\begin{verbatim}
arch_baseline_pin/parameter_sweeps/comsol_pscdeg_sweep_001.csv
\end{verbatim}
and is processed by
\begin{verbatim}
python .\scripts\plot_comsol_parameter_sweep.py
\end{verbatim}
The script writes:
\begin{itemize}
\item \path{outputs/tables/comsol_parameter_sweep_metrics.csv};
\item \path{outputs/tables/comsol_parameter_sweep_single_axis_metrics.csv};
\item combined J--V and metric figures under
\path{outputs/figures/comsol_parameter_sweep/};
\item the combined J--V examples figure and combined metric-grid figure; and
\item individual PDF, PNG, and SVG panels for the $N_{t,0}$, $R_{s,0}$, and
$R_{sh,0}$ sweeps in the same directory.
\end{itemize}
This sweep is a baseline electronic-parameter sensitivity study, not an aging
or lifetime simulation.

\section{Architecture-Stress Data Products}

The baseline architecture aging data are stored under the architecture folder
\path{arch_baseline_pin/Aging/}. The canonical light-temperature stress-matrix
J--V data file is:
\begin{verbatim}
Stress_matrix_jv_003.csv
\end{verbatim}
and is processed by:
\begin{verbatim}
python .\scripts\process_comsol_architecture_stress_jv.py
\end{verbatim}
The resulting tables and figures use the prefix
\path{baseline_pin_stress_matrix_jv_003_}. Superseded development data are
kept only in the repository audit trail and are not used for current
quantitative claims.

The internal-state profile grid is stored in:
\begin{verbatim}
arch_baseline_pin/Aging/profile_grid_light_temperature/
\end{verbatim}
It covers 300--400 K and 0.01--1.00 sun, with canonical filenames and an
original-file manifest. The grid is processed with:
\begin{verbatim}
python .\scripts\process_comsol_profile_grid.py
\end{verbatim}
and writes the long-format state summary under
the architecture-stress processed-data folder.
